\section{经由格点微扰论实现格点与连续统之间的匹配}

在本节中，我们旨在建立把格点上的准夸克分布与连续统中通常夸克分布联系起来的单圈匹配。在格点上，除图~\ref{1loopFeyn} 中的各图之外，还存在一幅额外的单圈图（示于图~\ref{1loopLQCD}）。图~\ref{1loopFeyn} 的第一幅图与这幅额外图不过是费曼规范下夸克场的波函数重正化，这在格点 QCD 中早已被充分理解；由夸克数守恒，它们的贡献也可以从图~\ref{1loopFeyn} 的其余图中得到。如文献~\cite{Xiong:2013bka} 与~\cite{Ji:2015jwa} 所示，在连续统中，图~\ref{1loopFeyn} 的第二、三幅图对强子动量只含对数依赖而非对数发散；其贡献与正规化无关，因而在连续统与格点上相同。人们只需把图~\ref{1loopFeyn} 最后一幅图的线性发散部分在格点与连续统之间进行匹配。这便给出格点截断与连续统单圈核中线性发散之间的一个联系~\cite{Xiong:2013bka}（一般而言二者会相差一个有限项）。为获得这一联系，我们基于一个简单的离散化规范作用量版本，用格点微扰论计算威尔逊线自能。


\begin{figure}[tbp]
\centering
\includegraphics[width=0.2\textwidth]{images/1loopL1}
\caption{格点上准夸克分布的额外单圈图，共轭图未画出。}
\label{1loopLQCD}
%\vspace*{-.5em}
\end{figure}


威尔逊线自能图的贡献可以写为
\begin{align}
\delta W&=\frac{\alpha_s C_F}{8\pi^3}\int d^4(k a_L) \frac{1-\cos(k_z z)}{(1-\cos(k_z a_L))4\sum_\lambda \sin^2\frac{k_\lambda a_L}{2}}\non\\
&=\alpha_s C_F[\frac{\pi}{2}\frac{|z|}{a_L}-\frac{1}{\pi}\ln\frac{\pi |z|}{a_L}-\frac{1}{\pi}]+\mathcal O(\frac{a_L}{z}),
\end{align}
其中 $a_L$ 表示格点间距。把它与式~\eqref{WLcoordspace} 的展开相比较，得到连续统与格点线性发散之间的如下匹配关系
\beq
\Lambda=\frac{\pi}{a_L}.
\eeq
用上述匹配条件替换文献~\cite{Xiong:2013bka} 单圈核中的 $\Lambda$，我们便得到联系格点上准 PDF 与通常 PDF 的单圈核。在式~\eqref{qUnpolDef} 的被积函数中加入因子 $e^{-\delta m |z|}$，即得不含幂次发散的改进准分布 $\tilde q_{\text{imp}}$
\beq\label{im}
   \tilde q_{\text{imp}}(x, \Lambda, p^z) = \int^\infty_{-\infty} \frac{dz}{4\pi} e^{izk^z-\delta m |z|}  \langle p|\overline{\psi}(0, 0_\perp, z)
   \gamma^z L(z,0)\psi(0) |p\rangle \ .
\eeq

改进准分布 $\tilde q_{\text{imp}}$ 与 PDF $q$ 之间的匹配核 $Z$ 定义为
\beq
\tilde q_{\text{imp}}(x, a_L, p^z)  = \int^1_{-1} \frac{dy}{|y|} Z\left(\frac{x}{y},{p^z a_L},\frac{\mu}{p^z}\right) q(y, \mu) +  {\cal O}\left(\Lambda^2_{\rm QCD}/(p^z)^2,  M^2/(p^z)^2\right)\ .
\eeq
单圈的 $Z$ 因子可以从文献~\cite{Xiong:2013bka} 的式 (15)--(19) 中提取，其中 $1/a_L$ 的贡献已被图~\ref{1loopFeynmassCT} 的反项图减除：
\begin{equation}
             Z\left(\xi,  p^z a_L,\frac{\mu}{p^z}\right) = \delta (\xi-1) + \frac{\alpha_s}{2\pi} Z^{(1)}
             \left(\xi, p^z a_L,\frac{\mu}{p^z}\right) + \dots ,
\end{equation}
其中
\begin{align}
Z^{(1)}/C_F
=\left\{ \begin{array} {ll} \left(\frac{1+\xi^2}{1-\xi}\right)\ln \frac{\xi}{\xi-1} + 1\ , & \xi>1\ , \\ \left(\frac{1+\xi^2}{1-\xi}\right)\ln \frac{(p^z)^2}{\mu^2}+\left(\frac{1+\xi^2}{1-\xi}\right)\ln\big[{4\xi(1-\xi)}\big]-\frac{2\xi}{1-\xi}+1\ , & 0<\xi<1\ , \\ \left(\frac{1+\xi^2}{1-\xi}\right)\ln \frac{\xi-1}{\xi}-1\ , & \xi<0\ , \end{array} \right.
\end{align}
且 $\xi=1$ 附近的额外贡献为
\begin{align}
\delta Z^{(1)}
=\frac{\alpha_S C_F}{2\pi} \int dy \left\{ \begin{array} {ll} -\frac{1+y^2}{1-y}\ln \frac{y}{y-1}-1\ , & y>1\ , \\ -\frac{1+y^2}{1-y}\ln \frac{(p^z)^2}{\mu^2}-\frac{1+y^2}{1-y}\ln\big[4y(1-y)\big]+\frac{2y(2y-1)}{1-y}+1\ , & 0<y<1\ , \\ -\frac{1+y^2}{1-y}\ln \frac{y-1}{y}+1\ , & y<0\ . \end{array} \right.
\end{align}
鉴于幂次发散已被消除，上述单圈匹配核有助于从格点数据可靠地提取通常的 PDF。
